% !TEX program = uplatex
\documentclass[uplatex,dvipdfmx,a4paper,11pt]{jsarticle}

\usepackage[margin=22mm]{geometry}
\usepackage{amsmath,amssymb,amsthm,mathtools}
\usepackage{mathpartir}
\usepackage{booktabs,array,tabularx}
\usepackage{enumitem}
\usepackage{xcolor}
\usepackage{hyperref}
\usepackage{pxjahyper}
\hypersetup{
  colorlinks=true,
  linkcolor=blue!48!black,
  citecolor=blue!48!black,
  urlcolor=blue!48!black,
  pdftitle={Weighted Effect Subtraction Core},
  pdfsubject={algorithmic specification for weighted effect subtraction},
  pdfauthor={}
}

\setlength{\parindent}{1em}
\setlength{\parskip}{0.22em}
\setlist[itemize]{topsep=0.35em,itemsep=0.15em,parsep=0pt}
\setlist[enumerate]{topsep=0.35em,itemsep=0.15em,parsep=0pt}
\allowdisplaybreaks

\newtheoremstyle{jpplain}
  {0.7em}{0.7em}{\normalfont}{}
  {\bfseries}{。}{0.5em}{}
\theoremstyle{jpplain}
\newtheorem{definition}{定義}[section]
\newtheorem{remark}[definition]{注意}
\newtheorem{invariant}[definition]{不変条件}
\newtheorem{example}[definition]{例}

\newcommand{\Fam}{\mathsf{Fam}}
\newcommand{\Id}{\mathsf{Id}}
\newcommand{\Weight}{\mathsf{Weight}}
\newcommand{\Ineq}{\mathsf{Ineq}}
\newcommand{\Set}{\mathsf{Set}}
\newcommand{\fin}{\mathsf{fin}}
\newcommand{\dom}{\mathsf{dom}}
\newcommand{\ub}{\mathsf{ub}}
\newcommand{\lb}{\mathsf{lb}}
\newcommand{\fresh}{\mathsf{fresh}}
\newcommand{\eff}{\mathbin{!}}
\newcommand{\row}[2]{[#1;#2]}
\newcommand{\emptyrow}{\varnothing}
\newcommand{\All}{\mathsf{All}}
\newcommand{\Empty}{\mathsf{Empty}}
\newcommand{\Unit}{\mathbf 1}
\newcommand{\subty}[3]{#1\mathrel{<:}_{#2,#3}}
\newcommand{\wocc}[2]{#1^{#2}}
\newcommand{\compose}{\mathbin{\oplus}}
\newcommand{\ann}{\mathsf{ann}}
\newcommand{\gen}{\mathsf{gen}}
\newcommand{\inst}{\mathsf{inst}}
\newcommand{\principal}{\mathsf{principal}}
\newcommand{\compact}{\mathsf{compact}}
\newcommand{\coalesce}{\mathsf{coalesce}}
\newcommand{\perform}{\mathsf{perform}}
\newcommand{\ret}{\mathsf{return}}
\newcommand{\handle}{\mathsf{handle}}
\newcommand{\with}{\mathsf{with}}
\newcommand{\seqeff}{\mathsf{seq}}

\title{Weighted Effect Subtraction Core\\
\large algorithmic core for weighted effect rows}
\author{}
\date{2026年6月21日}

\begin{document}
\maketitle

\begin{abstract}
この文書は weighted effect subtraction の algorithmic core を定義する。
重みの合成は結果重みと family 不等式を返す。
型比較は左右の重みを持つ judgment
\[
  \subty{S}{L}{R}\,T
\]
で扱う。
制約機械は subtype 制約、family 不等式、変数 bound、row subtraction、principal extraction からなる。
\end{abstract}

\tableofcontents

\section{表面項}

\begin{definition}[skeleton と表面項]
\[
\begin{array}{rcl}
 q &::=& \emptyrow \mid \ast \mid H,\\
 K &::=& q \mid K\to K,\\[0.3em]
 e &::=& () \mid \mathsf{ref}\ x \mid \lambda(x:K):q.\ e \mid e\ e \mid \perform_h(e)\\
   && \mid\ \handle\ e\ \with\ h(y,k)\Rightarrow e_h\mid \ret\,x\Rightarrow e_r\\
   && \mid\ \mathsf{let}\ x=e_1\ \mathsf{in}\ e_2\\
   && \mid\ \mathsf{let\ rec}\ f:K=e_1\ \mathsf{in}\ e_2.
\end{array}
\]
束縛名を値として取り出す表面操作は \(\mathsf{ref}\ x\) である。
\(\mathsf{ref}\) 出現は weighted occurrence を作る。
\end{definition}

\section{型と重み}

\subsection{型}

値型、computation 型、effect row を次で置く。
\[
\begin{array}{rcl}
 A &::=& \alpha \mid \Unit \mid C\to C,\\
 C &::=& A\eff R,\\
 R &::=& \rho \mid \emptyrow \mid \row H R.
\end{array}
\]
\(\row H R\) は head \(H\) と tail \(R\) を持つ effect row である。
weighted occurrence を \(\wocc{T}{W}\) と書く。

\subsection{weight}

\begin{definition}[family 集合]
\(\Fam\) は effect family atom の集合である。
family 集合の meet は集合交差である。
\end{definition}

\begin{definition}[weight]
\[
  \Weight = \Id \to_{\fin} (\Fam,\mathbb Z,\Fam).
\]
\(W(i)=(U,n,O)\) と書く。
第一成分 \(U\) は upper-bound family、第二成分 \(n\) は stack length、第三成分 \(O\) は own family である。
\[
  i\notin\dom(W)\quad\Rightarrow\quad W(i)=(\All,0,\All).
\]
\end{definition}

\begin{definition}[制約付き合成]
weight 合成は次の部分関数である。
\[
  \compose : \Weight\times\Weight \to \Weight\times \mathcal P_{\fin}(\Fam\times\Fam).
\]
\[
  W_1\compose W_2=(W_3,I)
\]
であるとき、各 \(i\in\dom(W_1)\cup\dom(W_2)\) について
\[
\begin{aligned}
 W_3(i)
  ={}&(\operatorname{pr}_1(W_1(i))\cap \operatorname{pr}_1(W_2(i)),\\
     &\operatorname{pr}_2(W_1(i))+\operatorname{pr}_2(W_2(i)),\\
     &\operatorname{pr}_3(W_1(i))\cap \operatorname{pr}_3(W_2(i))),
\end{aligned}
\]
また
\[
  I=\{\,(\operatorname{pr}_3(W_1(i)),\operatorname{pr}_1(W_2(i)))\mid
       i\in\dom(W_1)\cup\dom(W_2)\,\}.
\]
\end{definition}

\begin{definition}[合成制約の展開]
合成は追加制約 \(I\) を返すため、weight だけの代数演算として扱わない。
制約集合への展開を
\[
  \mathsf{ineq}(I)=\{\,L<:H\mid (L,H)\in I\,\}
\]
と書く。
\end{definition}

\section{weighted subtyping}

型比較の judgment は
\[
  \subty{S}{L}{R}\,T
\]
である。\(L\) は左側の occurrence までに蓄積した weight、\(R\) は右側の occurrence までに蓄積した weight である。

\subsection{occurrence weight}

weighted occurrence に当たったとき、左右それぞれの weight へ吸収する。
\[
\inferrule*[right=W-Occ]
 {W_1\compose L=(L_1,I_1)\\
  W_2\compose R=(R_2,I_2)\\
  \subty{S}{L_1}{R_2}\,T\\
 (lo<:hi)\ \text{for all }(lo,hi)\in I_1\cup I_2}
 {\subty{\wocc{S}{W_1}}{L}{R}\,\wocc{T}{W_2}}
\]
右側の occurrence は \(W_2\compose R\) の順で吸収する。

\subsection{構造分解}

関数型は domain で左右の weight を交換する。
\[
\inferrule*[right=W-Fun]
 {\subty{D_1}{R}{L}\,C_1\\
  \subty{C_2}{L}{R}\,D_2}
 {\subty{C_1\to C_2}{L}{R}\,D_1\to D_2}
\]

computation 型は値型と row へ分解する。
\[
\inferrule*[right=W-Comp]
 {\subty{A}{L}{R}\,B\\
  \subty{S}{L}{R}\,T}
 {\subty{A\eff S}{L}{R}\,B\eff T}
\]

この単純言語の型構造は、関数型と computation 型だけである。

\section{制約機械}

\subsection{制約と変数境界}

\[
\begin{array}{rcl}
c\in\mathsf{Con} &::=& \subty{S}{L}{R}\,T \mid F<:G,\\
\mathcal G,\mathcal P &\in& \mathcal P_{\fin}(\mathsf{Con}).
\end{array}
\]
\(\mathcal G\) は生成済み制約集合、\(\mathcal P\) は処理済み制約集合である。
worklist は \(\mathcal G\setminus\mathcal P\) である。

型変数状態 \(\sigma\) は weighted lower/upper bounds を持つ。
\[
\begin{aligned}
 \lb_\sigma(\alpha)&\subseteq \mathsf{Type}\times\Weight\times\Weight,\\
 \ub_\sigma(\alpha)&\subseteq \mathsf{Type}\times\Weight\times\Weight.
\end{aligned}
\]
\(\subty{T}{L}{R}\,\alpha\) は \(\alpha\) の lower bound として \((T,L,R)\) を保存する。
\(\subty{\alpha}{L}{R}\,T\) は \(\alpha\) の upper bound として \((T,L,R)\) を保存する。
状態遷移は \(\langle\sigma;\mathcal G;\mathcal P\rangle
\to\langle\sigma';\mathcal G';\mathcal P'\rangle\) と書く。
全ての追加は集合和である。

\subsection{variable-bound propagation}

lower/upper pair の weight 合成を次で定義する。
\[
\begin{array}{rcl}
\mathsf{bcomp}((L_s,R_s),(L_t,R_t))
 &=&
 ((L,R),I_L\cup I_R)\\[0.2em]
\end{array}
\]
ただし
\[
  L_s\compose L_t=(L,I_L),
  \qquad
  R_t\compose R_s=(R,I_R).
\]

bound pair から生成される制約集合:
\[
\mathsf{cgen}(S,L_s,R_s,T,L_t,R_t)
=
\{\subty{S}{L}{R}\,T\}\cup\mathsf{ineq}(I)
\]
ただし \(\mathsf{bcomp}((L_s,R_s),(L_t,R_t))=((L,R),I)\).

一つの lower bound 追加が生成する制約集合:
\[
\mathcal B^-_\sigma(\alpha,S,L_s,R_s)=
\bigcup_{\mathclap{(T,L_t,R_t)\in\ub_\sigma(\alpha)}}
\mathsf{cgen}(S,L_s,R_s,T,L_t,R_t)
\]

一つの upper bound 追加が生成する制約集合:
\[
\mathcal B^+_\sigma(\alpha,T,L_t,R_t)=
\bigcup_{\mathclap{(S,L_s,R_s)\in\lb_\sigma(\alpha)}}
\mathsf{cgen}(S,L_s,R_s,T,L_t,R_t)
\]

lower bound 挿入:
\[
\mathsf{ins}^-_\alpha(S,L,R)(\sigma,\mathcal G)=
\begin{cases}
(\sigma,\mathcal G),
  & (S,L,R)\in\lb_\sigma(\alpha),\\
(\sigma[\lb_\alpha:=\lb_\sigma(\alpha)\cup\{(S,L,R)\}],
 \mathcal G\cup\mathcal B^-_\sigma(\alpha,S,L,R)),
  & \text{otherwise}.
\end{cases}
\]

upper bound 挿入:
\[
\mathsf{ins}^+_\alpha(T,L,R)(\sigma,\mathcal G)=
\begin{cases}
(\sigma,\mathcal G),
  & (T,L,R)\in\ub_\sigma(\alpha),\\
(\sigma[\ub_\alpha:=\ub_\sigma(\alpha)\cup\{(T,L,R)\}],
 \mathcal G\cup\mathcal B^+_\sigma(\alpha,T,L,R)),
  & \text{otherwise}.
\end{cases}
\]

制約処理:
\[
\begin{array}{rcl}
\mathsf{step}(F<:G)(\sigma,\mathcal G)
  &=& (\sigma,\mathcal G)\quad(F\subseteq G),\\[0.3em]
\mathsf{step}(\subty{S}{L}{R}\,\alpha)(\sigma,\mathcal G)
  &=& \mathsf{ins}^-_\alpha(S,L,R)(\sigma,\mathcal G)
      \quad(S\notin\mathsf{Var}),\\[0.3em]
\mathsf{step}(\subty{\alpha}{L}{R}\,T)(\sigma,\mathcal G)
  &=& \mathsf{ins}^+_\alpha(T,L,R)(\sigma,\mathcal G)
      \quad(T\notin\mathsf{Var}),\\[0.3em]
\mathsf{step}(\subty{\alpha}{L}{R}\,\beta)(\sigma,\mathcal G)
  &=&
  \text{let }(\sigma_1,\mathcal G_1)=
      \mathsf{ins}^-_\beta(\alpha,L,R)(\sigma,\mathcal G)\\
  && \text{in }\mathsf{ins}^+_\alpha(\beta,L,R)(\sigma_1,\mathcal G_1).
\end{array}
\]

\[
\inferrule*[right=C-Step]
 {c\in\mathcal G\setminus\mathcal P\\
  \mathsf{step}(c)(\sigma,\mathcal G)=(\sigma',\mathcal G')}
 {\langle\sigma;\mathcal G;\mathcal P\rangle
  \to
  \langle\sigma';\mathcal G';\mathcal P\cup\{c\}\rangle}
\]

\begin{invariant}[variable-bound closure]
\(\mathcal G\setminus\mathcal P=\emptyset\) なら、全ての
\((S,L_s,R_s)\in\lb_\sigma(\alpha)\) と
\((T,L_t,R_t)\in\ub_\sigma(\alpha)\) について、
\[
  \mathsf{cgen}(S,L_s,R_s,T,L_t,R_t)\subseteq\mathcal G.
\]
\end{invariant}

\section{effect row subtraction}

\subsection{row upper への到達}

\begin{definition}[row upper]
row upper は
\[
  \subty{\alpha}{L}{R}\,\row{S}{\beta}
\]
である。
\end{definition}

\begin{enumerate}
  \item \(L=R=0\) のとき、
  \[
    \mathcal G\cup\{\subty{\alpha}{0}{0}\,\row{S}{\beta}\}
  \]
  を作る。

  \item それ以外では、まず
  \[
    L\compose R=(W,I)
  \]
  を作り、\((lo,hi)\in I\) から \(lo<:hi\) を生成する。

  \item \(W\) に残る own family と floor 情報から \(\mathsf{common}(W)\) を計算し、
  \[
    U=S\cap \mathsf{common}(W)
  \]
  を得る。

  \item \(U\neq\Empty\) なら fresh row 変数 \(\gamma\) を取り、
  \[
    \subty{\alpha}{0}{0}\,\row{U}{\gamma}
  \]
  を追加する。さらに subtraction
  \[
    W-U=(W',I')
  \]
  を計算し、\((lo,hi)\in I'\) から \(lo<:hi\) を生成したうえで、
  \[
    \subty{\gamma}{W'}{0}\,\beta
  \]
  を追加する。

  \item \(U=\Empty\) なら、残差変数を作らず
  \[
    \subty{\alpha}{W}{0}\,\beta
  \]
  を追加する。
\end{enumerate}

\subsection{residual slot}

\[
  \mathsf{residual}(\alpha,U,W')=\gamma
\]
は部分関数である。
同じ \((\alpha,U,W')\) には同じ \(\gamma\) を返す。
\(\beta\) は key に含めない。

\begin{invariant}[no self-fueling residual]
\(\subty{\gamma}{L}{R}\,\row{S}{\gamma}\) は新しい境界情報を持たない。
\end{invariant}

\section{注釈分解}

型注釈は極性に応じて lower 側と upper 側へ分解する。
判断を
\[
  \mathsf{decomp}^{p}(T)\leadsto (T^-,T^+,\mathcal C)
\]
と書く。\(p\in\{+,-,0\}\) は現在の極性である。

\[
\begin{array}{c|c|c}
p & q & \mathsf{decomp}^{p}(q)\\
\hline
+ & \ast & (\rho,\rho,\emptyset)\\
+ & H & (\rho,\rho,\{\rho<:\row{H}{\rho'}\})\\
+ & \emptyrow & (\rho,\rho,\{\rho<:\emptyrow\})\\
- & \ast & (\rho,\rho,\emptyset)\\
- & H & (\rho,\wocc{\rho}{\{a\mapsto(H,0,H)\}},\emptyset)\\
- & \emptyrow & (\rho,\wocc{\rho}{\{a\mapsto(\Empty,0,\Empty)\}},\emptyset)\\
0 & H & (\row{H}{\rho},\wocc{\rho}{\{a\mapsto(H,0,H)\}},\emptyset)
\end{array}
\]
\(\rho,\rho'\) と \(a\) は fresh である。
表面省略は \(q=\ast\) へ脱糖する。

\section{式の型付け}

\begin{definition}[型付け judgment]
\[
  \Gamma;\sigma \vdash e \Rightarrow C \dashv \sigma';\mathcal C.
\]
\end{definition}

\subsection{変数参照}

\[
\inferrule*[right=I-Ref]
 {\Gamma(x)=\sigma_x\\
  \inst(\sigma_x)=C}
 {\Gamma;\sigma\vdash \mathsf{ref}\ x\Rightarrow C\dashv\sigma;\emptyset}
\]
\(\inst\) は scheme 内の型変数、row 変数、weight id を freshen する。

\subsection{lambda}

\[
\inferrule*[right=I-Lam]
 {\mathsf{arg}(K)=C_x\\
  \Gamma,x:C_x;\sigma\vdash e\Rightarrow C_o\dashv\sigma';\mathcal C\\
  \ann(q,C_o)\leadsto \mathcal C_q}
 {\Gamma;\sigma\vdash \lambda(x:K):q.e
   \Rightarrow (C_x\to C_o)\eff\emptyrow
   \dashv\sigma';\mathcal C\cup\mathcal C_q}
\]
\(\ann(q,C_o)\) は返り effect slot に対する \(\mathsf{decomp}^{+}\) の制約である。

\subsection{application}

\[
\inferrule*[right=I-App]
 {\Gamma;\sigma\vdash e_1\Rightarrow A_1\eff E_1\dashv\sigma_1;\mathcal C_1\\
  \Gamma;\sigma_1\vdash e_2\Rightarrow A_2\eff E_2\dashv\sigma_2;\mathcal C_2\\
  B,\rho_i,\rho_o\ \fresh\\
  \subty{A_1}{0}{0}\,((A_2\eff\rho_i)\to(B\eff\rho_o))}
 {\Gamma;\sigma\vdash e_1\,e_2
   \Rightarrow B\eff\seqeff(E_1,E_2,\rho_o)
   \dashv\sigma_2;\mathcal C_1\cup\mathcal C_2}
\]

\subsection{let}

\[
\inferrule*[right=I-Let]
 {\Gamma;\sigma\vdash e_1\Rightarrow C_1\dashv\sigma_1;\mathcal C_1\\
  \principal_{\sigma_1}(C_1)=\Sigma_1\\
  \gen(\Gamma,\Sigma_1)=\sigma_x\\
  \Gamma,x:\sigma_x;\sigma_1\vdash e_2\Rightarrow C_2\dashv\sigma_2;\mathcal C_2}
 {\Gamma;\sigma\vdash \mathsf{let}\ x=e_1\ \mathsf{in}\ e_2
   \Rightarrow C_2\dashv\sigma_2;\mathcal C_1\cup\mathcal C_2}
\]

\begin{remark}
\(\principal_{\sigma}(C)\) は bounds graph から compact type を作る操作である。
ここで得た scheme を環境へ入れるため、let は単なる構文上の代入ではない。
\end{remark}

\section{principal extraction}

\[
  \principal_{\sigma}(C)
  =\gen(\compact(\coalesce_{\sigma}(C))).
\]
weight 付き変数出現を compact 化するとき、次の制限を置く。

\begin{invariant}[二次的な変数を開かない]
変数 \(\alpha\) を展開して現れた変数 \(\beta\) について、\(\beta\) の bounds をその場でさらに展開しない。
展開元の変数に直接保存されていた bounds だけを compact 化の対象にする。
\end{invariant}

\section{不変条件}

\begin{enumerate}
  \item weight 合成は、結果 weight と追加 family 不等式を同時に返す。
  \item 左右 weight は関数 domain で交換される。
  \item row subtraction に到達するまで、左右 weight を合成しない。
  \item residual slot は \((\alpha,U,W')\) で共有する。
  \item principal extraction では二次的な変数を開かない。
\end{enumerate}

\end{document}
